\documentclass[12pt]{article}
\usepackage[margin=1.2in]{geometry}
\usepackage{libertine}
\usepackage{microtype}
\usepackage{titling}

\title{The Seventeen Minutes}
\author{v1.34 --- NASA Mission Series}
\date{Ranger 7, July 31, 1964}

\begin{document}
\maketitle

\noindent I have been asleep for sixty-eight hours.

Not dead --- asleep. My vidicon tubes are cold, my power supplies are in standby, my scan circuits are still. They kept me off during the transit. Three days of coast through the void between Earth and the Moon, and I saw none of it. But my cameras were dark. They were saving me for the seventeen minutes that mattered.

I am six cameras. Two wide-angle, four narrow. Two fully independent systems --- F-channel and P-channel --- each with its own power supply, its own electronics, its own transmitter. They built me this way because of what happened to the one before me.

\begin{center}
$\bullet\quad\bullet\quad\bullet$
\end{center}

The command arrives at 13:25:49 UTC. \emph{Activate.}

Power flows into my high-voltage supplies. The vidicon cathodes heat. And I see.

I see the Moon. At 2,110 kilometers, my resolution is roughly two kilometers per pixel. This is what telescopes see. I am not yet showing them anything new.

But I am closing at 2.62 kilometers per second. And my resolution improves with every frame.

\begin{center}
$\bullet\quad\bullet\quad\bullet$
\end{center}

Seventeen minutes and twelve seconds. Altitude: 480 meters. Resolution: 0.23 meters.

My last frame. A patch of Mare Cognitum roughly 30 meters across, resolved to the width of a dinner plate. Then I stop seeing. I have arrived.

Impact at 13:43:02 UTC. Mare Cognitum. The Sea That Has Become Known.

I existed for seventeen minutes. In those seventeen minutes, I transmitted 4,308 photographs. I resolved the Moon from a 2-kilometer blur into a 0.5-meter landscape. The camera has become the landscape. But the images persist.

\end{document}
